% ==============================================================================
% Fichier : 07_evaluation_finale.tex
% Description : Evaluation Finale du Cours d'Audit des SI
% ==============================================================================

\chapter{Évaluation Finale du Cours}\label{chap:08_evaluation_finale}

Cette évaluation valide l'acquisition des concepts fondamentaux de l'audit, la maîtrise de la démarche méthodologique et la connaissance des référentiels.

% ------------------------------------------------------------------------------
\section{Partie A : QCM (Choix Multiples)}
% ------------------------------------------------------------------------------

\textbf{Consigne :} Cochez la ou les bonnes réponses.

\subsection*{Chapitre 2 : Fondamentaux}
\begin{enumerate}
    \item L'audit a pour sujet principal :
    \begin{itemize} \item[A.] Une opération unitaire. \item[B.] Un processus global. \item[C.] Un employé précis. \item[D.] Une transaction financière. \end{itemize}
    \item Le "Risque Résiduel" est :
    \begin{itemize} \item[A.] Le risque avant toute mesure de sécurité. \item[B.] Le risque qui reste après l'application des contrôles. \item[C.] Un risque qui ne peut pas se produire. \item[D.] Le risque de l'auditeur de se tromper. \end{itemize}
    \item La règle des "Quatre Yeux" signifie :
    \begin{itemize} \item[A.] Avoir quatre yeux (mutation génétique). \item[B.] Qu'une opération sensible doit être validée par deux personnes distinctes. \item[C.] Relire un document quatre fois. \item[D.] Avoir deux superviseurs. \end{itemize}
    \item Un contrôle "Préventif" a pour but de :
    \begin{itemize} \item[A.] Détecter l'erreur après coup. \item[B.] Empêcher l'erreur ou la fraude de se produire. \item[C.] Coriger l'erreur. \item[D.] Sanctionner le coupable. \end{itemize}
    \item L'assurance fournie par un audit est dite "Raisonnable" car :
    \begin{itemize} \item[A.] L'auditeur est gentil. \item[B.] Elle n'est pas absolue (limites de l'échantillonnage et du temps). \item[C.] Elle est garantie à 100\%. \item[D.] Elle dépend du prix payé. \end{itemize}
\end{enumerate}

\subsection*{Chapitre 3 : Démarche et Normes}
\begin{enumerate}
    \setcounter{enumi}{5}
    \item La norme qui encadre la méthodologie de l'audit est :
    \begin{itemize} \item[A.] ISO 27001. \item[B.] ISO 19011. \item[C.] COBIT. \item[D.] COSO. \end{itemize}
    \item Le principe de "Confidentialité" pour l'auditeur signifie :
    \begin{itemize} \item[A.] Ne rien dire à personne. \item[B.] Protéger les informations sensibles découvertes. \item[C.] Cacher les résultats à la direction. \item[D.] Signer un accord de non-divulgation. \end{itemize}
    \item Le "Rapport d'Orientation" est validé par :
    \begin{itemize} \item[A.] L'auditeur. \item[B.] L'audité. \item[C.] Le Commanditaire (PDG/COMEX). \item[D.] Le fournisseur informatique. \end{itemize}
    \item La phase de "Contradictoire" permet à l'audité de :
    \begin{itemize} \item[A.] Contester les observations et apporter des preuves nouvelles. \item[B.] Refuser le rapport. \item[C.] Corriger lui-même ses erreurs. \item[D.] Changer l'auditeur. \end{itemize}
    \item Une observation d'audit contient obligatoirement :
    \begin{itemize} \item[A.] Le nom du coupable. \item[B.] Le risque, la cause et le constat. \item[C.] Une amende. \item[D.] Une opinion personnelle. \end{itemize}
\end{enumerate}

\subsection*{Chapitre 4 : Référentiels}
\begin{enumerate}
    \setcounter{enumi}{10}
    \item COSO est un modèle de :
    \begin{itemize} \item[A.] Sécurité réseau. \item[B.] Contrôle Interne. \item[C.] Développement logiciel. \item[D.] Gestion de projet. \end{itemize}
    \item COBIT est un référentiel pour :
    \begin{itemize} \item[A.] La gouvernance et le management des SI. \item[B.] La comptabilité. \item[C.] La sécurité physique. \item[D.] Les ressources humaines. \end{itemize}
    \item La norme certifiable pour la sécurité de l'information est :
    \begin{itemize} \item[A.] ISO 27002. \item[B.] ISO 27001. \item[C.] ISO 27005. \item[D.] ISO 19011. \end{itemize}
    \item Le domaine COBIT "Acquérir et Implémenter" concerne :
    \begin{itemize} \item[A.] L'exploitation quotidienne. \item[B.] Les projets et changements. \item[C.] La stratégie. \item[D.] Le budget. \end{itemize}
    \item Les 5 composantes de COSO incluent :
    \begin{itemize} \item[A.] L'environnement de contrôle. \item[B.] La gestion des stocks. \item[C.] Le marketing. \item[D.] La paie. \end{itemize}
\end{enumerate}

\subsection*{Chapitre 5 : Types d'Audits}
\begin{enumerate}
    \setcounter{enumi}{15}
    \item L'audit d'une application en service vérifie :
    \begin{itemize} \item[A.] La fiabilité et la sécurité des données. \item[B.] Le projet de développement. \item[C.] Les contrats des fournisseurs. \item[D.] L'organigramme de l'entreprise. \end{itemize}
    \item Les "Contrôles Généraux" (General Controls) portent sur :
    \begin{itemize} \item[A.] Une fonction précise du logiciel. \item[B.] L'environnement global (OS, Réseau, Sauvegarde). \item[C.] Une facture spécifique. \item[D.] Un utilisateur unique. \end{itemize}
    \item L'audit de sécurité vérifie le principe DIC/T. La lettre I signifie :
    \begin{itemize} \item[A.] Informatique. \item[B.] Intégrité. \item[C.] Intelligence. \item[D.] Investissement. \end{itemize}
    \item L'audit de projet intervient généralement :
    \begin{itemize} \item[A.] Uniquement après la panne. \item[B.] Pendant le projet (jalons) ou à la fin. \item[C.] Avant le cahier des charges. \item[D.] Jamais. \end{itemize}
    \item Un "Pentest" (Test d'intrusion) est :
    \begin{itemize} \item[A.] Un audit de configuration. \item[B.] Une simulation d'attaque réelle. \item[C.] Une réunion de travail. \item[D.] Un test de performance. \end{itemize}
\end{enumerate}

\subsection*{Chapitre 6 : Sujets Connexes}
\begin{enumerate}
    \setcounter{enumi}{20}
    \item Le risque principal de l'utilisation de l'IA pour l'auditeur est :
    \begin{itemize} \item[A.] L'automatisation. \item[B.] Le biais algorithmique et la "boîte noire". \item[C.] Le coût élevé. \item[D.] La rapidité. \end{itemize}
    \item La fraude "interne" est commise par :
    \begin{itemize} \item[A.] Un concurrent. \item[B.] Un employé de l'entreprise. \item[C.] Un client. \item[D.] Un hacker étranger. \end{itemize}
    \item La "Séparation des Tâches" est un contrôle :
    \begin{itemize} \item[A.] Correctif. \item[B.] Détectif. \item[C.] Préventif. \item[D.] Technique. \end{itemize}
    \item Pour détecter une fraude, l'auditeur utilise souvent :
    \begin{itemize} \item[A.] L'intuition. \item[B.] L'analyse de données (Data Mining). \item[C.] Les rumeurs. \item[D.] La voyance. \end{itemize}
    \item La ligne éthique (Whistleblowing) sert à :
    \begin{itemize} \item[A.] Signaler anonymement des comportements suspects. \item[B.] Commander du matériel. \item[C.] Gérer les mots de passe. \item[D.] Contacter le support. \end{itemize}
\end{enumerate}

% ------------------------------------------------------------------------------
\section{Partie B : Vrai ou Faux}
% ------------------------------------------------------------------------------

\textbf{Consigne :} Justifiez si faux (optionnel).

\begin{enumerate}
    \item L'auditeur interne est obligatoirement un employé de l'entreprise. \textbf{(Vrai)}
    \item Un contrôle détectif empêche la fraude de se produire. \textbf{(Faux - Il la détecte après coup)}
    \item L'indépendance est un principe éthique obligatoire pour l'auditeur. \textbf{(Vrai)}
    \item Le rapport d'audit est la seule production de l'auditeur. \textbf{(Faux - Il y a les papiers de travail, le rapport d'orientation...)}
    \item COSO est un référentiel réservé à l'informatique. \textbf{(Faux - C'est un modèle de contrôle interne généraliste)}
    \item ISO 27002 contient les exigences de certification. \textbf{(Faux - C'est ISO 27001, l'ISO 27002 est un guide de bonnes pratiques)}
    \item L'audit d'application vérifie uniquement le code source. \textbf{(Faux - Il vérifie les données, les processus et l'environnement)}
    \item Le "Suivi d'audit" (Follow-up) sert à vérifier l'application des recommandations. \textbf{(Vrai)}
    \item L'IA peut être utilisée comme outil d'aide à l'audit. \textbf{(Vrai)}
    \item Le principe de "Présentation équitable" permet à l'auditeur de mentir pour protéger l'entreprise. \textbf{(Faux - Il doit être honnête et factuel)}
\end{enumerate}

% ------------------------------------------------------------------------------
\section{Partie C : Questions Ouvertes}
% ------------------------------------------------------------------------------

\textbf{Consigne :} Répondez en 10-15 lignes.

\begin{enumerate}
    \item Expliquez la différence fondamentale entre l'Approche fondée sur les risques et l'Approche fondée sur les contrôles en audit.
    \item Quels sont les 5 composantes du référentiel COSO et quel est leur lien ?
    \item Décrivez les 3 phases principales de la relation avec l'audité lors d'une mission (Ouverture, Terrain, Clôture).
    \item Pourquoi la "Charte d'Audit" est-elle un document indispensable pour l'auditeur interne ?
    \item Comparez les objectifs d'un "Audit de Sécurité" et d'un "Audit de Performance".
    \item Quels sont les risques spécifiques liés à l'intégration de l'Intelligence Artificielle dans les SI que l'auditeur doit vérifier ?
    \item Expliquez le principe de séparation des tâches (Règle des deux mains) avec un exemple concret.
\end{enumerate}

% ------------------------------------------------------------------------------
\section{Partie D : Cas Pratique (Étude de Cas)}
% ------------------------------------------------------------------------------

\subsection{Scénario : Audit de l'Application "Comptabilité" de Lizbiz's}

\begin{displayquote}
Vous êtes l'auditeur senior mandaté pour auditer le module comptabilité de l'APPLICATION de Lizbiz's.
Lors de vos tests, vous constatez les faits suivants :
\begin{enumerate}
    \item L'administrateur système "Root" a accès directement à la base de données financières sans passer par l'application.
    \item Les sauvegardes journalières sont stockées sur un disque dur USB posé sur le bureau du comptable.
    \item Le comptable principal part en vacances et donne son mot de passe à son stagiaire pour qu'il continue le travail.
\end{enumerate}
\end{displayquote}

\textbf{Travaux demandés :}

\begin{enumerate}
    \item Pour chaque fait constaté, identifiez :
    \begin{itemize}
        \item Le risque encouru par Lizbiz's.
        \item Le type de contrôle défaillant (Préventif, Détectif, organisationnel, technique).
    \end{itemize}
    \item Rédigez une recommandation appropriée pour le constat n°3 (Partage de mot de passe).
    \item En vous référant à ISO 27001, citez un domaine de l'Annexe A qui aurait permis d'éviter le constat n°2 (Sauvegardes).
\end{enumerate}
